%% !TEX root = manuscript.tex

\chapter{La reprise des données vers SPO : des opérations aux tiers}
\minitoc
\newpage

Ce chapitre expose le chantier le plus lourd de l'année : la reconstitution, dans le nouvel
ERP, du socle de données du groupe. Il s'est déroulé en deux temps. Le premier semestre a
porté sur les structures de gestion, opérations, tranches et budgets, migrées depuis la base
historique Grimmo. Le second a porté sur les tiers, c'est-à-dire les clients acquéreurs,
repris depuis le CRM Vadimm. Les deux missions partagent une méthode et se distinguent par la
question qu'elles posent : injecter des structures suppose de comprendre un modèle de données
inconnu~; injecter des personnes suppose de décider ce qui fait qu'une personne est la même.

\hypertarget{ce-que-le-changement-derp-imposait}{%
\section{Ce que le changement d'ERP imposait}\label{ce-que-le-changement-derp-imposait}}

Le remplacement de Grimmo par SPO a commandé le calendrier de toute l'année, et trois
contraintes en découlaient. Une contrainte de \textbf{volumétrie} d'abord : 283 opérations devaient
être reprises avec leurs subdivisions et leurs budgets, échelle à laquelle une saisie manuelle
intégrale était exclue. Une contrainte de \textbf{continuité de service} ensuite : les utilisateurs
devaient pouvoir facturer dès le basculement, une interruption de deux semaines étant déjà
inacceptable. Une contrainte de \textbf{conformité au modèle cible} enfin : SPO impose ses propres
règles d'intégrité, qui n'étaient pas toutes documentées et qu'il a fallu découvrir en les
heurtant.

La Figure 1.1 situe ces deux reprises dans le système d'information du groupe. Les trois
outils et leurs usages métier occupent le haut du schéma, reliés par les flèches qui figurent
la circulation normale de la donnée~; les tableurs forment la bande du bas, et les traits
pointillés qui les rattachent à chacun d'eux marquent les échanges qui ne passent par aucun
système. On y lit le problème que les deux missions de ce chapitre attaquent : le même client
pouvait exister dans le CRM et dans l'ERP sans qu'aucun lien ne les rattache.

\begin{figure}[!htbp]
\centering
\includegraphics[width=0.65\textwidth]{fig_1_1_systeme_information.png}
\caption{Les outils du système d'information et la circulation de la donnée}
\end{figure}

À ces contraintes s'ajoute une perspective absente du calendrier initial : le passage prochain
vers l'infrastructure technique du groupe Nexity, qui a influé sur certains choix techniques,
notamment lorsqu'il a fallu évaluer si l'architecture retenue pour le projet de Data Science
tiendrait à une échelle supérieure.

\hypertarget{la-migration-des-opuxe9rations}{%
\section{La migration des opérations}\label{la-migration-des-opuxe9rations}}

Cette mission, conduite au premier semestre, est exposée en détail dans mon mémoire de
mi-parcours. Je n'en reprends ici que ce dont la suite a besoin : la méthode, les deux
contraintes que l'ERP impose, et l'état dans lequel le socle se trouvait à la fin du semestre.

\hypertarget{muxe9thode-et-industrialisation-des-flux}{%
\subsection{Méthode et industrialisation des flux}\label{muxe9thode-et-industrialisation-des-flux}}

Face aux 283 opérations du périmètre, une saisie manuelle intégrale était exclue. Nous avons
retenu une stratégie hybride : injection automatisée des gros volumes structurels via les API,
traitement manuel des cas historiques ne pouvant entrer dans les gabarits. Vouloir tout
automatiser aurait conduit à écrire des règles pour des cas uniques, ce qui coûte plus cher
que de les traiter à la main.

J'ai d'abord analysé la documentation technique de l'API de Salvia {[}2{]}, puis testé
manuellement les \emph{endpoints} pour décortiquer la mécanique de création d'une donnée. Nous
avons ensuite déployé le module DirectAPI, dont j'ai construit le pipeline d'alimentation sous
la forme de maquettes Excel en trois étapes : \textbf{extraction} de la base source par MyReport
Builder, \textbf{transformation et mapping} par formules RechercheX, \textbf{chargement} via DirectAPI.
La transformation concentre l'essentiel du travail, la nomenclature de SPO étant stricte : un
libellé de poste budgétaire ou un type de tranche non reconnu fait rejeter la ligne. Excel a
été retenu parce que la maquette reste relisible par les personnes qui connaissent le métier
et non le code.

Le processus a été sécurisé par une validation itérative : 3 opérations pilotes, puis un lot
de 12, avant la migration de masse. Cette progression par paliers rend les erreurs lisibles,
et c'est elle que j'ai reconduite au second semestre.

\hypertarget{deux-contraintes-de-lerp-et-la-documentation-qui-en-est-sortie}{%
\subsection{Deux contraintes de l'ERP, et la documentation qui en est sortie}\label{deux-contraintes-de-lerp-et-la-documentation-qui-en-est-sortie}}

La première est l'ordre strict d'injection, contrainte d'intégrité référentielle que la
documentation de l'éditeur n'explicitait pas. Je l'ai découverte en analysant les codes
d'erreur renvoyés par l'API : il faut créer l'opération, puis les tranches de travaux, puis
les tranches commerciales, puis le budget. La Figure 1.2 représente cette séquence~; toute
injection qui ne la respecte pas échoue, le système ne trouvant pas l'objet parent.

\begin{figure}[!htbp]
\centering
\includegraphics[width=0.65\textwidth]{fig_1_2_sequence_injection.png}
\caption{La séquence d'injection imposée par l'intégrité référentielle de l'ERP}
\end{figure}

J'ai alors rédigé une documentation interne cartographiant cette séquence et les prérequis de
chaque objet, parce que ces règles n'existaient nulle part ailleurs que dans mes notes de
test. C'est le premier livrable écrit de mon année qui n'était pas un fichier de données, et
c'est celui qui a le plus servi aux autres : la reprise des tiers, exposée plus bas, s'est
appuyée directement dessus.

La seconde tient aux identifiants. Invisibles dans l'interface, ils changent à chaque mise à
jour : l'API ne remplace pas un budget, elle clôture la version précédente et en crée une
nouvelle. Ne pouvant me fier à aucun identifiant statique, j'ai conçu un requêtage dynamique
en trois temps --- un appel GET pour lister les versions, un tri pour retenir la plus récente,
un appel POST sur cet identifiant. J'en retire une leçon plus générale : une API ne se
contente pas d'exposer des données, elle impose un modèle de comportement, et comprendre
l'intention derrière la contrainte est ce qui permet de la contourner proprement.

\hypertarget{le-duxe9ploiement-en-coquilles-vides-et-le-bilan-du-semestre}{%
\subsection{Le déploiement en «~coquilles vides~» et le bilan du semestre}\label{le-duxe9ploiement-en-coquilles-vides-et-le-bilan-du-semestre}}

Les utilisateurs devaient facturer dès le basculement, notamment pour le lien avec la Gestion
Électronique des Documents (GED). Nous avons donc adopté une stratégie de déploiement en
«~coquilles vides~», qui dissocie le contenant --- la structure de l'opération --- du contenu,
c'est-à-dire les données budgétaires détaillées, en injectant d'abord le code opération et le
libellé.

C'est la décision dont j'ai le plus appris en conduite de projet. Elle suppose d'accepter
qu'un système soit temporairement incomplet et de le documenter comme tel : une donnée
manquante que personne ne signale finit par être considérée comme absente pour de bon. Tenir
la liste de ce qui restait à reprendre a transformé une dette implicite en dette suivie.

\begin{longtable}[]{@{}
  >{\raggedright\arraybackslash}p{(\columnwidth - 2\tabcolsep) * \real{0.4000}}
  >{\raggedright\arraybackslash}p{(\columnwidth - 2\tabcolsep) * \real{0.6000}}@{}}
\toprule\noalign{}
\begin{minipage}[b]{\linewidth}\raggedright
Objet du modèle SPO
\end{minipage} & \begin{minipage}[b]{\linewidth}\raggedright
État à la fin du premier semestre
\end{minipage} \\
\midrule\noalign{}
\endhead
\bottomrule\noalign{}
\endlastfoot
Opérations & Créées pour l'ensemble du périmètre, 283 dossiers \\
Tranches de travaux & Créées \\
Tranches commerciales & Création en cours \\
Budgets détaillés & Reportés, stratégie des «~coquilles vides~» \\
Historique des factures & Non repris \\
\end{longtable}

Sur un chantier de reprise, l'avancement ne se mesure pas au nombre d'enregistrements créés
mais à la part du modèle cible réellement couverte : créer 283 opérations sans leurs budgets,
ce n'est pas avoir fait la moitié du travail, c'est avoir rendu la seconde moitié possible.

\hypertarget{la-reprise-des-tiers}{%
\section{La reprise des tiers}\label{la-reprise-des-tiers}}

Conduite au second semestre, cette mission prolonge directement la précédente et constitue le
premier pas vers l'interopérabilité entre l'ERP et le CRM : elle transfère le référentiel
client, la synchronisation permanente sans double saisie restant à construire. Elle change
toutefois de nature : il ne s'agit plus d'injecter des structures de gestion, mais des
personnes, avec ce que cela implique de risques de doublons et d'exigence de qualité.

\hypertarget{contexte-et-enjeux-de-la-reprise-des-tiers}{%
\subsection{Contexte et enjeux de la reprise des tiers}\label{contexte-et-enjeux-de-la-reprise-des-tiers}}

\hypertarget{du-crm-vadimm-uxe0-lerp-spo-le-puxe9rimuxe8tre-des-donnuxe9es-clients}{%
\subsubsection{Du CRM Vadimm à l'ERP SPO : le périmètre des données clients}\label{du-crm-vadimm-uxe0-lerp-spo-le-puxe9rimuxe8tre-des-donnuxe9es-clients}}

Après la création des opérations et de leurs tranches, il restait à reprendre les données
commerciales, et en premier lieu les clients. Dans SPO, ceux-ci sont désignés par le terme
générique de \emph{tiers} : la même entité de base représente un acquéreur particulier comme une
société, et se voit attribuer un ou plusieurs rôles. Reprendre les tiers, c'était alimenter le
socle sur lequel toute la gestion commerciale allait s'appuyer.

Le périmètre était d'importer l'ensemble des clients ayant déjà acheté chez Angelotti. La
matière première venait du CRM Vadimm, sous la forme de deux extractions complémentaires :
celle des ventes enregistrées, qui porte 1~557 lignes, et celle des coordonnées des acquéreurs
et co-acquéreurs, qui en porte 1~602. Les deux ne se recouvrent donc pas exactement~; l'écart
de 45 lignes n'est pas documenté dans les fichiers sources et s'explique vraisemblablement par
des contacts renseignés sans vente correspondante. J'ai retenu le plus large des deux totaux
pour ne perdre aucun acquéreur : le périmètre de travail compte \textbf{1~602 lignes}.

La difficulté ne tenait pas au volume mais à la nature de la donnée. Une ligne de l'extraction
ne décrit pas un client, elle décrit une vente : un acquéreur qui a acheté trois lots apparaît
trois fois, et rien dans le fichier ne le signale. Transformer un fichier de ventes en
référentiel de personnes suppose de reconstituer l'identité de chacune, ce que le système
source ne faisait pas à ma place.

\hypertarget{lenjeu-central-cruxe9er-sans-dupliquer}{%
\subsubsection{L'enjeu central : créer sans dupliquer}\label{lenjeu-central-cruxe9er-sans-dupliquer}}

L'enjeu se résume à une contrainte simple à énoncer et difficile à tenir : créer chaque client
une fois et une seule. Deux sources de duplication se cumulaient, et elles n'appelaient pas la
même parade.

La première est interne au fichier : chaque ligne correspondant à une vente, un client fidèle
y figure autant de fois qu'il a acheté. Importer le fichier tel quel aurait créé autant de
fiches que d'achats, rendant impossible de reconstituer son historique ou de lui adresser un
courrier unique. La seconde est externe : au moment de préparer l'import, SPO n'était pas vide,
environ 13~000 tiers y ayant déjà été créés par saisie ou par d'autres canaux. Injecter les
1~602 lignes sans les confronter à cet existant aurait produit des doublons plus insidieux,
opposant une fiche importée et une fiche saisie, toutes deux légitimes en apparence.

Une reprise manuelle était donc exclue : vérifier ligne à ligne qu'un acquéreur n'existe ni
ailleurs dans le fichier ni déjà dans un référentiel de 13~000 tiers représentait un volume de
contrôles incompatible avec les délais, et un risque d'erreur que rien n'aurait permis de
tracer. J'ai construit un fichier de travail qui automatise ces contrôles par formules, de
sorte que la décision d'importer une ligne soit calculée et non prise au cas par cas.

\hypertarget{conception-du-fichier-de-travail-une-chauxeene-de-traitement-pilotuxe9e-par-formules}{%
\subsection{Conception du fichier de travail : une chaîne de traitement pilotée par formules}\label{conception-du-fichier-de-travail-une-chauxeene-de-traitement-pilotuxe9e-par-formules}}

\hypertarget{une-architecture-en-trois-temps-sources-feuille-pivot-feuilles-de-sortie}{%
\subsubsection{Une architecture en trois temps : sources, feuille pivot, feuilles de sortie}\label{une-architecture-en-trois-temps-sources-feuille-pivot-feuilles-de-sortie}}

Le fichier ne se présente pas comme un tableur de saisie mais comme une chaîne de traitement,
dans la lignée des maquettes d'injection du premier semestre : les données entrent d'un côté
sous leur forme brute et ressortent de l'autre sous la forme attendue par SPO, sans qu'aucune
donnée soit recopiée à la main. La Figure 1.3 en présente les trois étages.

\begin{enumerate}
\def\labelenumi{\arabic{enumi}.}

\item
  \textbf{Les feuilles sources}, en lecture seule : \texttt{VADIMM} porte les 1~557 lignes de ventes,
  \texttt{Builder\ coordonnée} les 1~602 lignes de coordonnées, et \texttt{SPO} l'extraction des tiers déjà
  existants, qui sert de référentiel de comparaison.
\item
  \textbf{La feuille pivot \texttt{BuilderTiers}}, cœur du dispositif : elle reçoit l'intégralité des
  lignes sources et y adjoint les colonnes de calcul qui déterminent, pour chacune, s'il
  s'agit d'un tiers nouveau, d'un doublon interne ou d'un tiers déjà présent dans SPO.
\item
  \textbf{Les six feuilles de sortie}, une par fichier d'import attendu par SPO, entièrement
  générées par formules et destinées à être exportées en CSV.
\end{enumerate}

\begin{figure}[!htbp]
\centering
\includegraphics[width=0.65\textwidth]{fig_1_3_chaine_traitement.png}
\caption{La chaîne de traitement à trois étages du fichier de travail}
\end{figure}

Une règle figure en tête de la note de procédure : on ne modifie jamais directement une
feuille de sortie. Ces feuilles étant calculées, toute correction saisie dedans serait écrasée
au recalcul suivant, sans laisser de trace. Toute intervention passe donc par \texttt{BuilderTiers}
ou par la feuille de coordonnées. C'est la contrepartie de l'automatisation : elle garantit
qu'à tout moment le contenu des fichiers d'import résulte des règles décrites et non d'une
retouche oubliée.

\hypertarget{le-moduxe8le-de-donnuxe9es-des-tiers-dans-spo-couples-personnes-physiques-personnes-morales}{%
\subsubsection{Le modèle de données des tiers dans SPO : couples, personnes physiques, personnes morales}\label{le-moduxe8le-de-donnuxe9es-des-tiers-dans-spo-couples-personnes-physiques-personnes-morales}}

Avant d'automatiser, il fallait comprendre comment SPO représente un tiers, par lecture de la
documentation puis tests unitaires de création. Deux caractéristiques ont structuré tout le
reste du travail.

La première est que SPO distingue trois \textbf{natures} de tiers, et que la nature commande la
manière dont l'identité est portée. Une \textbf{personne morale} est une société, identifiée par sa
raison sociale et son SIRET. Une \textbf{personne physique} est un individu, identifié par nom,
prénom et date de naissance. Un \textbf{couple} est une entité particulière : le tiers n'est ni l'un
ni l'autre des conjoints, il les regroupe, et les deux personnes qui le composent sont
enregistrées comme ses \emph{représentants}. Cette nature est de loin la plus fréquente, l'essentiel
de nos ventes se faisant à des couples acquéreurs.

La seconde est qu'un tiers n'est pas un enregistrement unique mais un ensemble d'objets liés :
la fiche principale ne porte que l'identité, tandis que l'adresse, le courriel, le téléphone,
les représentants et les rôles sont autant d'objets rattachés, avec chacun son fichier
d'import. C'est ce qui explique que mon fichier produise six feuilles de sortie et non une
seule. Cela explique aussi une asymétrie décisive : adresses, courriels et téléphones se
rattachent directement à une personne physique ou morale, mais pour un couple ils se rattachent
aux représentants. Cette différence, anodine sur le papier, s'est révélée être la principale
contrainte technique de l'import.

\hypertarget{les-trois-muxe9canismes-algorithmiques}{%
\subsection{Les trois mécanismes algorithmiques}\label{les-trois-muxe9canismes-algorithmiques}}

La feuille pivot repose sur trois mécanismes qui s'enchaînent : identifier les doublons
internes, écarter les tiers déjà créés dans SPO, puis distribuer les lignes retenues vers les
six feuilles de sortie. Ce sont eux qui distinguent un fichier de mise en forme d'un outil de
reprise, puisqu'ils déplacent la décision d'importer une ligne de l'opérateur vers une règle
écrite.

\hypertarget{la-duxe9duplication-par-cluxe9s-dunicituxe9-diffuxe9renciuxe9es-selon-la-nature-du-tiers}{%
\subsubsection{La déduplication par clés d'unicité différenciées selon la nature du tiers}\label{la-duxe9duplication-par-cluxe9s-dunicituxe9-diffuxe9renciuxe9es-selon-la-nature-du-tiers}}

Dédupliquer suppose de répondre à une question qui n'est pas technique : à partir de quoi
décide-t-on que deux lignes désignent la même personne~? C'est le point sur lequel j'ai le
plus hésité, et celui qui commande la qualité du résultat.

Une clé unique appliquée à tout le fichier ne pouvait pas convenir, précisément parce que les
trois natures de tiers ne s'identifient pas de la même façon : deux sociétés peuvent porter des
raisons sociales très proches sans être la même entité, alors qu'il est très improbable que
deux lignes partageant nom, prénom, date de naissance et adresse désignent deux personnes
distinctes. J'ai donc retenu \textbf{trois clés d'unicité distinctes}, une par nature, calculées
chacune dans sa propre colonne :

\begin{enumerate}
\def\labelenumi{\arabic{enumi}.}

\item
  \textbf{Pour un couple} : nom et prénom de l'acquéreur principal, désigné ADV dans nos
  extractions, complétés du nom et du prénom du co-acquéreur, désigné CO. C'est la
  combinaison des deux identités qui fait l'unicité, et non chacune prise isolément.
\item
  \textbf{Pour une personne physique} : le nom, le prénom, la date de naissance et l'adresse. La
  date de naissance sépare les homonymes, l'adresse ajoute un second point de comparaison.
\item
  \textbf{Pour une personne morale} : la raison sociale et le SIRET, identifiant légal de
  l'établissement, la raison sociale le doublant lorsqu'il fait défaut.
\end{enumerate}

Chaque colonne calcule un \textbf{rang} : la première occurrence d'un tiers reçoit le rang 1, les
suivantes les rangs 2, 3, et ainsi de suite. Seules les lignes de rang 1 sont retenues, les
autres correspondant à des achats successifs d'un client déjà compté. Le rang a un avantage
sur une suppression : aucune ligne n'est effacée, les doublons restent visibles, et l'on peut
à tout moment vérifier pourquoi une ligne n'a pas été importée.

Ce choix de trois clés différenciées est le véritable cœur méthodologique du travail, parce
qu'il traduit une idée simple : la règle d'unicité n'est pas une propriété du fichier, c'est
une propriété du métier. Chercher une clé universelle aurait conduit soit à fusionner des
tiers distincts, soit à en dupliquer, et dans les deux cas l'erreur n'aurait été découverte
que bien plus tard, dans l'usage quotidien de l'ERP.

\hypertarget{lexclusion-des-tiers-duxe9juxe0-pruxe9sents-dans-spo}{%
\subsubsection{L'exclusion des tiers déjà présents dans SPO}\label{lexclusion-des-tiers-duxe9juxe0-pruxe9sents-dans-spo}}

Pour chaque ligne de la feuille pivot, trois formules de recherche interrogent l'extraction
des tiers déjà créés dans SPO : les \textbf{personnes morales} par leur raison sociale, les
\textbf{personnes physiques} par leur nom et leur prénom, et les \textbf{couples} par les quatre noms
qui les composent, \textbf{en testant les deux ordres possibles} entre l'acquéreur et le
co-acquéreur.

Ce dernier point illustre le genre de détail qui décide de la réussite d'une reprise. Rien ne
garantit que l'ordre des deux conjoints soit le même dans Vadimm et dans SPO~; une recherche à
sens unique laisserait donc passer une partie des doublons, et les tiers concernés seraient
créés une seconde fois sans que personne s'en aperçoive avant longtemps.

Lorsqu'une de ces formules renvoie un code tiers de SPO, du type \texttt{T00012666}, la ligne est
exclue de l'import. \textbf{40 tiers} ont ainsi été détectés comme déjà existants. Rapporté aux
1~474 tiers uniques, le volume peut sembler faible, mais chacun de ces 40 cas aurait produit
un doublon durable dans le référentiel client.

La Figure 1.4 récapitule l'effet cumulé de ces deux mécanismes : des 1~602 lignes de départ,
la déduplication interne en écarte 128 qui correspondent à des achats répétés, ramenant le
total à 1~474 tiers uniques~; l'exclusion en retire 40 de plus, pour \textbf{1~434 tiers à créer}.
Chaque étage correspond à une règle explicite, écrite dans une colonne, donc vérifiable et
reproductible.

\begin{figure}[!htbp]
\centering
\includegraphics[width=0.65\textwidth]{fig_1_4_entonnoir_volumetrie.png}
\caption{De 1~602 lignes de départ à 1~434 tiers créés : l'effet des deux filtres successifs}
\end{figure}

\hypertarget{la-propagation-vers-les-feuilles-de-sortie-par-numuxe9rotation-suxe9quentielle}{%
\subsubsection{La propagation vers les feuilles de sortie par numérotation séquentielle}\label{la-propagation-vers-les-feuilles-de-sortie-par-numuxe9rotation-suxe9quentielle}}

Les six feuilles de sortie n'ont ni la même population ni le même nombre de lignes : celle des
tiers reçoit un enregistrement par tiers créé, celle des adresses ne concerne que les
personnes physiques et morales, celle des courriels seulement celles dont le courriel est
renseigné, celle des représentants deux lignes par couple plus une par société. Il fallait
donc à chaque feuille une numérotation propre.

Le mécanisme retenu attribue à chaque ligne de la feuille pivot un \textbf{numéro séquentiel par
destination} : une ligne sans vocation à alimenter une feuille ne reçoit pas de numéro pour
elle, les autres sont numérotées de 1 à \emph{n}. Chaque feuille de sortie n'a plus qu'à chercher,
pour sa ligne \emph{i}, celle de \texttt{BuilderTiers} portant le numéro \emph{i} dans la colonne qui la
concerne.

Ce détour rend les six feuilles indépendantes tout en les gardant synchronisées avec une
source unique : si une ligne est corrigée ou si une exclusion apparaît, les six numérotations
se recalculent et les fichiers restent cohérents. C'est exactement le problème qui, en saisie
manuelle, produit les incohérences les plus difficiles à détecter --- une adresse rattachée au
mauvais tiers parce qu'une ligne a été insérée sans que le reste suive.

La volumétrie des six feuilles traduit directement le modèle de données décrit en 1.3.2, et le
tableau de la section 1.3.6 la récapitule : 1~434 lignes pour la liste maîtresse des tiers et
autant pour le rôle «~Client~», mais 1~871 représentants puisque chaque couple en compte deux,
et seulement 553 adresses, 538 téléphones et 469 courriels, ces trois fichiers ne portant que
les personnes physiques et morales.

\begin{quote}
\textbf{À insérer --- figure.} Capture d'écran de la feuille \texttt{BuilderTiers}, montrant côte à
côte les colonnes de données reprises et les colonnes calculées : les trois rangs de
déduplication avec des valeurs 1 et 2 visibles, et une ligne portant un code SPO dans les
colonnes de rapprochement. Anonymiser les noms d'acquéreurs. Les figures suivantes sont à
renuméroter après insertion.
\end{quote}

\hypertarget{de-la-pruxe9paration-uxe0-linjection-procuxe9dure-et-contraintes-techniques}{%
\subsection{De la préparation à l'injection : procédure et contraintes techniques}\label{de-la-pruxe9paration-uxe0-linjection-procuxe9dure-et-contraintes-techniques}}

\hypertarget{guxe9nuxe9ration-et-normalisation-des-fichiers-dimport}{%
\subsubsection{Génération et normalisation des fichiers d'import}\label{guxe9nuxe9ration-et-normalisation-des-fichiers-dimport}}

Transformer les six feuilles en fichiers réellement acceptés par SPO a mobilisé quatre
contraintes de format que la documentation ne mentionnait pas et que les premiers rejets ont
fait apparaître : un encodage UTF-8 avec le point-virgule pour séparateur~; la conversion en
ASCII pur de tous les noms, prénoms, adresses et professions, SPO refusant les caractères
accentués~; la suppression de la seconde ligne d'en-tête~; et le découpage en lots de 200
tiers au maximum, homogènes en nature. Cette dernière règle explique que l'import des couples
se soit fait en cinq lots successifs, puisqu'un même fichier ne peut mélanger les natures. Le
détail figure en annexe~A.

Ces contraintes sont consignées dans la note de procédure remise à l'équipe, avec la marche à
suivre pour rejouer l'import s'il devait être différé : d'autres tiers pouvant être créés
entre la préparation et l'injection, la procédure prévoit de redemander une extraction
actualisée, de remplacer la feuille de référence et de recalculer, les tiers devenus existants
sortant automatiquement du périmètre. Le fichier reste ainsi utilisable dans le temps, ce qui
était l'un de mes objectifs de conception.

\hypertarget{le-contournement-par-api-872-requuxeates-patch-pour-les-coordonnuxe9es-des-couples}{%
\subsubsection{Le contournement par API : 872 requêtes PATCH pour les coordonnées des couples}\label{le-contournement-par-api-872-requuxeates-patch-pour-les-coordonnuxe9es-des-couples}}

C'est ici qu'est apparue la contrainte la plus structurante du projet. Le format d'import CSV
ne permet pas de porter les courriels et les téléphones des couples : pour un couple, ces
coordonnées ne sont pas des attributs du tiers mais de ses représentants, et les fichiers
d'import de coordonnées ne savent viser qu'un tiers. Sur les 881 couples créés, les CSV
pouvaient donc produire l'identité et les représentants, mais pas les coordonnées de contact ---
plus de six tiers sur dix seraient arrivés dans l'ERP sans adresse électronique ni téléphone.

Saisir ces coordonnées à la main représentait plusieurs centaines de saisies et reconduisait
le risque d'erreur que le fichier automatisé visait à supprimer. J'ai donc retenu la seconde
voie, l'API de SPO, déjà pratiquée au premier semestre. Elle expose une opération de
modification partielle au format \textbf{JSON Patch}, normalisé par la RFC~6902 {[}3{]}, qui décrit une
modification non pas en renvoyant l'objet entier mais sous la forme d'opérations élémentaires
indiquant l'action, le chemin de la propriété visée et la valeur à appliquer. C'était
exactement ce qu'il fallait : ajouter un courriel et un téléphone sur un représentant déjà
créé, sans toucher au reste de sa fiche.

J'ai construit le corps de chaque requête à partir de la feuille \texttt{Builder\ coordonnée}, puis
regroupé l'ensemble dans une collection Postman de \textbf{872 requêtes PATCH}, organisées en cinq
lots correspondant aux cinq fichiers d'import de couples. Les cinq lots ont été exécutés après
la création des couples, et l'intégralité des requêtes est passée sans erreur.

\begin{quote}
\textbf{À insérer --- figure.} Capture d'écran de la collection Postman : arborescence des cinq
lots et vue d'une requête PATCH sur son onglet \emph{Body}, montrant la structure JSON Patch.
Masquer l'adresse du serveur et les identifiants. Les figures suivantes sont à renuméroter
après insertion.
\end{quote}

Ce contournement m'a appris quelque chose que je n'avais pas formulé aussi clairement : dans
un projet de reprise, la limite d'un format d'import n'est pas nécessairement une limite du
système. Le fichier plat et l'API sont deux portes d'entrée vers le même modèle de données, et
savoir basculer de l'une à l'autre au bon moment, plutôt que d'accepter une perte de données
ou de retomber dans la saisie manuelle, est une compétence que la reprise du premier semestre
n'avait pas eu l'occasion de mobiliser.

\hypertarget{traitement-des-cas-particuliers-et-arbitrages}{%
\subsection{Traitement des cas particuliers et arbitrages}\label{traitement-des-cas-particuliers-et-arbitrages}}

Une chaîne pilotée par formules ne traite pas tout. Neuf tiers ont demandé une intervention
manuelle, répartis en quatre familles : trois adresses hors de France, que SPO refusait de
créer avec un pays étranger et qui ont été créées en France puis rectifiées~; trois tiers dont
la nature était mal renseignée dans Vadimm, corrigée dans la feuille pivot avant génération
puisque la nature commande la clé d'unicité et la feuille de sortie~; un tiers relevant d'une
opération de dation, déjà présent dans SPO~; et deux doublons portant deux adresses
divergentes. L'annexe~B donne le détail de chaque cas.

Ce dernier est le seul qui ait demandé un arbitrage véritable : la déduplication savait dire
qu'il s'agissait de la même personne, elle ne pouvait pas dire où cette personne habite
aujourd'hui. Après vérification des dates dans Vadimm, nous avons conservé l'adresse rattachée
à la vente la plus récente.

Ces neuf cas délimitent précisément ce que l'automatisation peut et ne peut pas faire : les
formules savent appliquer une règle, elles ne savent pas décider laquelle de deux adresses
également plausibles est la bonne. La valeur du fichier n'est pas d'avoir supprimé la décision
humaine, mais de l'avoir concentrée sur une poignée de cas identifiés au lieu de la diluer sur
1~602 lignes.

\hypertarget{bilan-de-limport-et-points-de-vigilance-pour-la-suite}{%
\subsection{Bilan de l'import et points de vigilance pour la suite}\label{bilan-de-limport-et-points-de-vigilance-pour-la-suite}}

L'intégralité des imports CSV et des requêtes PATCH a été exécutée avec succès.

\begin{longtable}[]{@{}
  >{\raggedright\arraybackslash}p{(\columnwidth - 4\tabcolsep) * \real{0.6364}}
  >{\raggedleft\arraybackslash}p{(\columnwidth - 4\tabcolsep) * \real{0.1212}}
  >{\raggedright\arraybackslash}p{(\columnwidth - 4\tabcolsep) * \real{0.2424}}@{}}
\toprule\noalign{}
\begin{minipage}[b]{\linewidth}\raggedright
Objet créé dans SPO
\end{minipage} & \begin{minipage}[b]{\linewidth}\raggedleft
Volume
\end{minipage} & \begin{minipage}[b]{\linewidth}\raggedright
Canal
\end{minipage} \\
\midrule\noalign{}
\endhead
\bottomrule\noalign{}
\endlastfoot
Tiers & 1~434 & Import CSV \\
Adresses & 553 & Import CSV \\
Courriels (personnes physiques et morales) & 469 & Import CSV \\
Téléphones (personnes physiques et morales) & 538 & Import CSV \\
Représentants (couples et sociétés) & 1~871 & Import CSV \\
Rôles «~Client~» & 1~434 & Import CSV \\
Coordonnées des couples & 872 & Requêtes PATCH \\
\end{longtable}

Les 1~434 tiers se répartissent en \textbf{881 couples}, \textbf{462 personnes physiques} et \textbf{91
personnes morales}, répartition que représente la Figure 1.5. Cette structure justifie a
posteriori l'attention portée aux couples, qui représentent plus de six tiers sur dix et
concentrent la contrainte de format qui a imposé le recours à l'API.

\begin{figure}[!htbp]
\centering
\includegraphics[width=0.65\textwidth]{fig_1_5_repartition_natures.png}
\caption{Répartition des 1~434 tiers créés selon leur nature}
\end{figure}

\begin{quote}
\textbf{À insérer --- figure.} Capture d'écran d'une fiche de couple dans SPO, onglet des
représentants ouvert, montrant les deux personnes du couple et les coordonnées ajoutées
par requête PATCH sur l'une d'elles. Anonymiser les noms.
\end{quote}

Quatre manipulations cassent le fichier sans produire le moindre message d'erreur : modifier
une feuille de sortie, modifier les colonnes calculées de la feuille pivot, insérer une ligne
au milieu d'une feuille, et trier la feuille pivot, dont la logique de rang est cumulative sur
l'ordre des lignes. La note de procédure les énonce avant même d'expliquer comment se servir
du fichier.

Les pièges de formule rencontrés relèvent tous du même défaut, et c'est l'enseignement
principal que j'en tire : dans une chaîne pilotée par formules, l'erreur ne se manifeste
presque jamais par un message d'erreur, mais par un compteur qui ne tombe pas juste. Chaîne
vide comptée comme une valeur, code de mois interprété comme des minutes, heures parasites
rendant différentes deux dates identiques, cellule vide confondue avec chaîne vide : chacun a
produit des compteurs surévalués ou des faux positifs, et aucun n'a signalé quoi que ce soit
(annexe~C). J'ai donc ajouté à la feuille pivot un panneau de compteurs que la procédure
impose de vérifier avant tout export : c'est le seul dispositif de détection dont je
disposais. Avec environ 60~000 cellules de formules, elle recommande enfin de travailler en
calcul manuel puis de forcer un recalcul complet avant chaque export.

\hypertarget{ce-que-ces-deux-reprises-ont-en-commun}{%
\section{Ce que ces deux reprises ont en commun}\label{ce-que-ces-deux-reprises-ont-en-commun}}

Les deux chantiers ont suivi la même méthode, que je n'avais pas choisie a priori : elle s'est
imposée parce que le modèle cible était inconnu. On ne peut pas spécifier complètement une
reprise dont on ignore les règles d'acceptation~; on avance donc par paliers de volume
croissant, en traitant chaque rejet comme une information sur le modèle cible plutôt que comme
une erreur à corriger. Sa propriété la plus utile est que \textbf{le coût d'une erreur y est
proportionnel au palier où on la découvre}~; sa limite est qu'elle démarre lentement et produit
d'abord une documentation interne plutôt que des données dans le système.

Les difficultés techniques se ramènent presque toutes à une même situation : un système
imposait une règle qui n'était écrite nulle part. L'ordre d'injection en est le cas le plus
net, le versionnement des budgets le second, le refus des caractères accentués et
l'impossibilité de porter les coordonnées des couples en CSV relèvent de la même famille.

Le régime de preuve de ces deux missions est celui de la \textbf{réconciliation} : la question n'est
pas si un modèle est performant, mais si ce qui est arrivé correspond à ce qui est parti. Elle
a pris la forme de la validation par paliers sur la migration et du panneau de compteurs sur
la reprise des tiers. Son angle mort mérite d'être nommé : elle prouve qu'une donnée est
arrivée, pas qu'elle est juste. Si une clé d'unicité était mal choisie, l'erreur passerait le
contrôle sans être vue, puisque les compteurs tomberaient juste.

Le référentiel commercial est ainsi reconstitué selon des règles écrites et rejouables. Le
chantier n'est pas clos pour autant : la reprise des historiques budgétaires reste devant
nous, et le fichier devra être rejoué avant la bascule définitive.
